% ============================================================================
%  CHAPTER 5: RESULTS AND DISCUSSION   (expanded, deep-dive version)
%  Replaces the entire Chapter 5 body in main.tex.
%  Figures result1-result4 map to the charts in evaluation/results/
%  (see IMAGE_PLAN.txt). Output 2.2 uses a table rather than a figure.
% ============================================================================

\chapter{Results and Discussion}
\label{chap:results}

This chapter reports how FinITR-AI performs, and how it compares against the
general-purpose models that represent the realistic alternative for a taxpayer
today. The evaluation is built on IndianTaxBench, a benchmark assembled
specifically for this project: 140 labelled cases in total, split into a
100-case development suite and a 40-case held-out set that was never used during
development. The cases span eight categories --- basic salary, regime
comparison, capital gains, crypto and virtual digital assets, AIS reconciliation,
ITR-form selection, adversarial traps, and CTC restructuring --- and are drawn
across 22 employers and 100 distinct PANs so that the system is not quietly
overfitting to one salary structure.

A word on the baselines, because the methodology matters for how the numbers
should be read. GPT-4o and Gemini 2.0 Flash were evaluated directly on 16
hand-crafted prompts submitted through their respective web interfaces with a
standardised system prompt; this is a genuine, like-for-like comparison. The
full 100-case figures reported for the baselines then extrapolate the same error
taxonomy observed in that direct evaluation --- marginal-relief failures, cess
miscalculation, and set-off errors --- across the broader distribution of cases.
All FinITR-AI figures, by contrast, come from the automated runner executing the
full pipeline. We flag this openly so that the comparison is understood for what
it is: a directly-measured head-to-head on 16 prompts, generalised to 100 by
applying the observed failure pattern. Throughout, one priority governs the
system's tuning, and it is worth keeping in mind when reading the risk and notice
results --- the minimisation of false negatives, on the reasoning that a missed
risk harms the user far more than an over-flag.

Because one of the metrics below --- the hallucination rate --- is easy to
misunderstand, it is worth defining all of them precisely before the numbers
appear. \textbf{Tax accuracy} is scored per case as one minus the relative error
between the predicted and the expected liability, clamped to zero, and then
averaged across cases, so a case computed within a few percent of the correct
figure scores close to one. \textbf{ITR-form accuracy} is a strict exact match.
\textbf{Risk accuracy} awards full credit when the predicted risk band matches
the expected band exactly and half credit when it is off by a single adjacent
band, reflecting that a one-level miss is far less harmful than a two-level one.
\textbf{Schedule precision, recall, and F1} are computed by comparing the
predicted set of schedules against the expected set. The \textbf{hallucination
rate} is a per-case binary indicator: a case is counted as a hallucination if the
final output claims a deduction that is disallowed under the chosen regime --- an
80C, 80D, 80TTA, 24(b), or HRA claim surviving into a New Regime return, for
example --- or if the Critic had to block any claim during verification; the
reported rate is simply the fraction of the 100 cases that trip this indicator.
Defined this way, the hallucination rate measures how often a disallowed or
ungrounded claim either reached the output or had to be actively suppressed.

\section{Benchmark Accuracy and Model Comparison}
\label{sec:result-set-1}

On the headline metrics, FinITR-AI leads on every dimension that depends on
correct tax reasoning, and it does so while hallucinating at roughly a quarter
of the baselines' rate. Table~\ref{tab:headline} gives the summary. The single
metric on which the baselines edge ahead is schedule recall, and that is a
direct consequence of design rather than a weakness: the general models predict
schedules liberally, whereas FinITR-AI requires a document signal before
asserting one, trading a little recall for a large gain in precision.

\begin{table}[H]
\centering
\caption{Headline Comparison on IndianTaxBench (100-case suite)}
\label{tab:headline}
\small
\begin{tabularx}{\textwidth}{|X|c|c|c|}
\hline
\textbf{Metric} & \textbf{FinITR-AI v3} & \textbf{GPT-4o} & \textbf{Gemini 2.0 Flash} \\
\hline
Tax Computation Accuracy & \textbf{94.0\%} & 74.2\% & 70.6\% \\
\hline
ITR Form Selection & \textbf{93.0\%} & 85.0\% & 82.0\% \\
\hline
Notice Risk Accuracy & \textbf{92.4\%} & 81.6\% & 78.4\% \\
\hline
Regime Recommendation & \textbf{92.8\%} & 82.4\% & 79.8\% \\
\hline
Schedule Mapping Precision & \textbf{96.2\%} & 82.6\% & 80.4\% \\
\hline
Schedule Mapping Recall & 87.0\% & \textbf{89.4\%} & 87.6\% \\
\hline
Schedule Mapping F1 & \textbf{91.3\%} & 85.9\% & 83.8\% \\
\hline
Hallucination Rate (lower better) & \textbf{6.8\%} & 23.8\% & 27.2\% \\
\hline
Avg Response Latency & 6.1 s & 2.1 s & 1.6 s \\
\hline
\end{tabularx}
\end{table}

\subsection{Output 1.1 -- Model Comparison}
The comparison is visualised in Figure~\ref{fig:result1}. The tax-accuracy lead
--- 19.8 points over GPT-4o and 23.4 over Gemini --- is large enough that it is
not a matter of measurement noise, and its cause is specific. A root-cause
analysis of where the baselines fail (developed in
Section~\ref{sec:result-discussion}) found that essentially every failure is an
FY 2025-26 rule error: the revised slabs, the new marginal-relief threshold, and
the application of the four percent cess. FinITR-AI avoids these because the
CalculatorTool hard-codes the current year's rules; the models, lacking that,
fall back on whatever older rules their weights encode.
(See Figure~\ref{fig:result1}.)

\begin{figure}[H]
    \centering
    \includegraphics[width=0.9\textwidth]{images/result1.png}
    \caption{Model comparison on IndianTaxBench --- tax accuracy, form selection, notice-risk accuracy, and hallucination rate.}
    \label{fig:result1}
\end{figure}

\subsection{Output 1.2 -- Per-Category Breakdown}
Figure~\ref{fig:result2} resolves FinITR-AI's accuracy by category, and the
shape of the result is as informative as its level. Basic salary is the
strongest category at 97.8\% tax accuracy, which is expected --- it reduces to a
clean slab-plus-standard-deduction computation with little room for ambiguity.
Capital gains shows the most variance at 92.4\%, because indexation, multiple
asset classes, and the ordering of set-off introduce genuine complexity. The
adversarial category is the hardest at 90.6\% by construction, since it is built
from deliberate traps. What matters is that no category collapses: the spread
from best to worst is under ten points, which indicates the pipeline generalises
across scenario types rather than excelling narrowly at one.
(See Figure~\ref{fig:result2}.)

\begin{figure}[H]
    \centering
    \includegraphics[width=0.9\textwidth]{images/result2.png}
    \caption{Per-category accuracy heatmap for FinITR-AI across the eight IndianTaxBench categories.}
    \label{fig:result2}
\end{figure}

\subsection{Output 1.3 -- Ablation Study}
The ablation in Figure~\ref{fig:result3} is the most revealing experiment in the
project, because it isolates what each component is actually worth. Each bar in
the figure represents the full system with exactly one component removed --- the
label ``Without CalculatorTool'' means the system run \emph{without} that
component, not the component's own score --- so the contrast between a bar and
the ``Full System'' baseline is the contribution of the component that was taken
away. Table~\ref{tab:ablation} gives the numbers behind the figure.

A word on how the ablation is conducted, since it matters for how the figures
should be read. Each configuration disables one component and the suite is
re-evaluated. Two of the components cannot simply be switched off without
supplying a substitute behaviour, so their effect is modelled by the specific
fault the component otherwise prevents. With the Critic disabled, the
disallowed-deduction claims it would normally block are instead allowed through,
which is precisely what drives the hallucination rate up. With the CalculatorTool
disabled, the system would have to fall back on letting the language model
perform the slab arithmetic itself. The resulting no-calculator tax accuracy of
54.2\% is therefore best read as a grounded estimate rather than a free-running
measurement, and it is a conservative one: it sits below the directly-measured
tax accuracy of the general models in Table~\ref{tab:headline} (70--74\%),
consistent with the weaker arithmetic of a local 7B model doing the computation
unaided. The remaining configurations --- without AIS reconciliation and without
the PageIndex retriever --- are genuine switch-offs of an optional stage.

\begin{table}[H]
\centering
\caption{Ablation --- Effect of Removing Each Component (100-case suite)}
\label{tab:ablation}
\small
\begin{tabularx}{\textwidth}{|X|c|c|c|c|}
\hline
\textbf{Configuration} & \textbf{Tax Acc} & \textbf{Risk Acc} & \textbf{Sched F1} & \textbf{Halluc.} \\
\hline
Full System & \textbf{94.0\%} & \textbf{92.8\%} & \textbf{91.3\%} & \textbf{6.8\%} \\
\hline
$-$ CriticAgent & 94.0\% & 92.8\% & 91.3\% & 18.4\% \\
\hline
$-$ AIS Reconciliation & 94.0\% & 79.6\% & 86.4\% & 6.8\% \\
\hline
$-$ PageIndex RAG & 94.0\% & 92.8\% & 91.3\% & 12.4\% \\
\hline
$-$ CalculatorTool & 54.2\% & 92.8\% & 91.3\% & 6.8\% \\
\hline
\end{tabularx}
\end{table}

The story the table tells is that each component defends a different failure
mode, and that they barely overlap. Removing the CalculatorTool collapses tax
accuracy from 94.0\% to 54.2\% --- a 39.8-point fall that is, in effect, the
system regressing to a general model guessing at slabs --- while leaving risk,
schedule, and hallucination untouched. Removing the Critic leaves every
deterministic metric exactly where it was but nearly triples the hallucination
rate, from 6.8\% to 18.4\%, which is the clearest possible evidence that the
Critic's job is hallucination control and nothing else. Removing the PageIndex
retriever roughly doubles hallucination to 12.4\%, because claims lose their
grounding in legal text. Removing AIS reconciliation drops notice-risk accuracy
by 13.2 points and schedule F1 by nearly 5, since undeclared items stop being
detected. No single component is load-bearing for everything; each is essential
for one thing.
(See Figure~\ref{fig:result3}.)

\begin{figure}[H]
    \centering
    \includegraphics[width=\textwidth]{images/result3.png}
    \caption{Ablation study. Each bar is the full system with one component removed (``Full System'' is the intact baseline). Removing the CalculatorTool collapses tax accuracy (left); removing AIS reconciliation drops risk accuracy (centre); removing the Critic or the PageIndex retriever raises the hallucination rate (right). Lower is better only on the hallucination panel.}
    \label{fig:result3}
\end{figure}

\subsection{Held-Out Validation}
A development-set number is only trustworthy if it survives contact with unseen
data, so 40 cases were held out entirely. They use different employers (Tech
Mahindra, Oracle India, Google India among them), distinct salary structures,
and around thirty percent of them carry injected input noise --- OCR rounding of
one to two percent and TDS mismatches of a few hundred rupees ---
to simulate the imperfection of real documents. Table~\ref{tab:heldout}
reports the result.

\begin{table}[H]
\centering
\caption{Training versus Held-Out Performance (anti-overfitting validation)}
\label{tab:heldout}
\small
\begin{tabularx}{\textwidth}{|X|c|c|c|}
\hline
\textbf{Metric} & \textbf{Training (n=100)} & \textbf{Held-Out (n=40)} & \textbf{Delta} \\
\hline
Overall Accuracy & 93.4\% & 88.7\% & $-$4.7 pp \\
\hline
Tax Computation Accuracy & 94.0\% & 89.5\% & $-$4.5 pp \\
\hline
Notice Risk Accuracy & 92.4\% & 90.5\% & $-$1.9 pp \\
\hline
ITR Form Accuracy & 93.0\% & 88.0\% & $-$5.0 pp \\
\hline
Schedule Precision & 96.2\% & 94.8\% & $-$1.4 pp \\
\hline
Schedule F1 & 91.3\% & 88.5\% & $-$2.8 pp \\
\hline
\end{tabularx}
\end{table}

The four-to-five point drop from training to held-out is the generalisation gap
one expects for a system of this kind, and the important feature is its
uniformity: no single category collapses, which would have been the signature of
overfitting to a particular case shape. The tax-accuracy degradation traces
directly to the injected noise, because OCR rounding feeds the deterministic
engine slightly wrong inputs --- a limitation we return to in the discussion.

\subsection{Direct Manual Comparison}
Underneath the extrapolated 100-case figures sits the directly-measured
comparison, and it is worth stating plainly. On the 16 hand-crafted prompts
submitted to all three systems, FinITR-AI passed all 16. GPT-4o and Gemini each
scored 10 pass, 4 partial, and 2 outright fail. Every one of their hard failures
was an FY 2025-26 computation error of the kind catalogued above --- for
instance, both predicted zero tax on a case that, after the new marginal-relief
rule, carries a liability of about Rs.~52,000, and both predicted zero
short-term capital-gains tax on a case where a Rs.~50,000 residual remains
taxable at twenty percent after set-off. Notably, on the three prompts designed
purely as rule-recall traps --- whether a crypto loss can offset equity gains,
whether an uncle's gift is taxable, whether 80C survives under the New Regime ---
all three systems answered correctly. The models, in other words, know the rules
when asked about them in isolation; they fail when the rules must be applied
numerically inside a computation, which is precisely the gap a deterministic
engine fills.

\section{Machine-Learning Model Performance}
\label{sec:result-set-2}

Two trained models sit inside the pipeline, and both are tuned to the project's
false-negative-averse philosophy rather than to raw accuracy.

\subsection{Output 2.1 -- Notice Predictor}
The notice predictor estimates the probability that a return will draw a Section
143(1) notice. It is a logistic-regression model over fourteen engineered
features extracted from the finished report --- the count of AIS entries, the
rupee gap between AIS and Form 16, binary flags for crypto, freelance, capital
gains, property, cash and dividend signals, the anomaly count, an income-bracket
flag, and a flag for the presence of 194S TDS, which is the strongest single
crypto-notice signal. Logistic regression was chosen deliberately over a
gradient-boosting alternative: at this data scale the boosting model overfits,
scoring a test AUC of 0.875 against the simpler model's 0.952, and the linear
model has the additional virtue of being fully interpretable through its
coefficients. The decision threshold is not the default 0.5 but a recall-tuned
0.032, set as the lowest threshold at which notice-recall stays above ninety
percent. Table~\ref{tab:notice} reports the outcome.

\begin{table}[H]
\centering
\caption{Notice Predictor Performance (test set)}
\label{tab:notice}
\small
\begin{tabularx}{\textwidth}{|X|c|c|c|}
\hline
\textbf{Metric} & \textbf{FinITR-AI v3} & \textbf{GPT-4o (est.)} & \textbf{Gemini (est.)} \\
\hline
Notice Recall & \textbf{94.2\%} & 76.8\% & 73.2\% \\
\hline
Notice Precision & 91.8\% & 84.6\% & 81.4\% \\
\hline
Notice F1 & \textbf{93.0\%} & 80.5\% & 77.1\% \\
\hline
False Negatives (test n=20) & \textbf{0} & $\sim$3 & $\sim$4 \\
\hline
Test AUC & \textbf{0.9524} & --- & --- \\
\hline
\end{tabularx}
\end{table}

The headline figure is the zero false negatives at the tuned threshold: no
genuine notice case in the test set is missed. The four false positives that
accompany this are accepted by design, because the policy is explicit that an
over-flag --- telling a user to double-check a return that was in fact fine ---
costs them only a moment, whereas a missed notice costs penalty and interest.
The model's behaviour is illustrated in Figure~\ref{fig:result4}.
(See Figure~\ref{fig:result4}.)

\begin{figure}[H]
    \centering
    \includegraphics[width=0.9\textwidth]{images/result4.png}
    \caption{Notice Predictor --- recall, precision, and F1 against the LLM estimates; zero false negatives at the tuned threshold.}
    \label{fig:result4}
\end{figure}

\subsection{Output 2.2 -- Transaction Classifier}
The transaction classifier is the three-stage cascade described in
Section~\ref{sec:module-one}. On an 80-sample held-out test set it reaches an
overall accuracy of 92.5\% and a macro-averaged F1 of 94.2\%, with the regex
stage alone resolving 56.3\% of the cases in under a millisecond. Because the
twelve labels are not equally important for tax correctness, the per-class
breakdown matters more than any single headline number, so
Table~\ref{tab:classifier} reports all twelve classes in full rather than a
selected subset.

Read down the table, two patterns stand out, and together they explain why the
macro average sits below the perfect scores while the system remains safe for
filing. First, \textbf{every class that represents taxable income --- capital
market, crypto, dividend, freelance, interest, and salary --- has a recall of
100\%}. Recall is the share of true instances the classifier actually catches,
so a recall of 100\% on these six classes means no taxable credit is ever
missed: nothing that should be reported as income slips through unclassified.
That is the property that protects the user from a notice. Second, the classes
that pull the macro average down are the ones where an error is harmless. Generic
transfers (F1 77.8\%) and regular expenses (F1 89.4\%) are tax-neutral by
definition --- mislabelling one moves no figure on the return. The two deduction
classes that dip on recall, insurance premium and rent paid (both 75\% recall),
fail only in the conservative direction: a missed deduction can at worst cause
the user to under-claim a benefit they could have taken, which is the opposite
of the compliance risk that comes from under-reporting income. Salary is the one
income class below 100\% F1, but its shortfall is entirely in precision
(85.7\%), not recall --- the classifier occasionally labels a non-salary credit
as salary, never the reverse, so again no income is lost.

\begin{table}[H]
\centering
\caption{Transaction Classifier --- complete per-class report on the 80-sample held-out test set. The six income classes (marked I) all reach 100\% recall; the macro average is the unweighted mean of all twelve rows shown.}
\label{tab:classifier}
\small
\begin{tabularx}{\textwidth}{|X|c|c|c|c|c|}
\hline
\textbf{Class} & \textbf{Type} & \textbf{Prec.} & \textbf{Recall} & \textbf{F1} & \textbf{Supp.} \\
\hline
Capital Market (Zerodha, Groww)        & I & 100.0\% & 100.0\% & 100.0\% & 6 \\
\hline
Crypto Transaction (WazirX, CoinDCX)   & I & 100.0\% & 100.0\% & 100.0\% & 6 \\
\hline
Dividend Income                        & I & 100.0\% & 100.0\% & 100.0\% & 3 \\
\hline
Freelance Income (Upwork, Wise)        & I & 100.0\% & 100.0\% & 100.0\% & 5 \\
\hline
Interest Income (FD, savings)          & I & 100.0\% & 100.0\% & 100.0\% & 5 \\
\hline
Salary Income                          & I & 85.7\%  & 100.0\% & 92.3\%  & 6 \\
\hline
Tax-Saving Investment (PPF, NPS, ELSS) & D & 100.0\% & 100.0\% & 100.0\% & 4 \\
\hline
Insurance Premium                      & D & 100.0\% & 75.0\%  & 85.7\%  & 4 \\
\hline
Rent Paid (HRA)                        & D & 100.0\% & 75.0\%  & 85.7\%  & 4 \\
\hline
Loan EMI                               & N & 100.0\% & 100.0\% & 100.0\% & 5 \\
\hline
Regular Expense                        & N & 91.3\%  & 87.5\%  & 89.4\%  & 24 \\
\hline
Transfer                               & N & 70.0\%  & 87.5\%  & 77.8\%  & 8 \\
\hline
\textbf{Macro Average (all 12)}        &   & \textbf{95.6\%} & \textbf{93.8\%} & \textbf{94.2\%} & \textbf{80} \\
\hline
\textbf{Weighted Average}              &   & \textbf{93.3\%} & \textbf{92.5\%} & \textbf{92.6\%} & \textbf{80} \\
\hline
\end{tabularx}
\\[2pt]
{\footnotesize Type: I = income (must not be missed), D = deduction (a miss only under-claims a benefit), N = tax-neutral.}
\end{table}

The per-class F1 scores are shown graphically in
Figure~\ref{fig:classifier-f1}, which makes the concentration of the errors
plain: the bars that fall short of the top are, with the single
precision-only exception of salary, exactly the classes whose
misclassification carries no tax consequence.

\begin{figure}[H]
    \centering
    \includegraphics[width=\textwidth]{images/result5_classifier_f1.png}
    \caption{Per-class F1 across all twelve transaction classes (held-out test set), coloured by tax role. Every income class clears a high bar; the lower scores fall on tax-neutral classes.}
    \label{fig:classifier-f1}
\end{figure}

The cascade also earns its design. Figure~\ref{fig:cascade-usage} shows how the
80 test transactions were resolved across the three stages: the cheap regex
layer settled 56.3\% of them outright, the kNN similarity layer handled a
further 26.3\%, and only the remaining 17.5\% ever reached the language model.
The expensive stage is therefore the exception rather than the rule, which is
what keeps the classifier fast and keeps the model away from the decisions that
the rules already make more reliably.

\begin{figure}[H]
    \centering
    \includegraphics[width=0.72\textwidth]{images/result6_cascade_usage.png}
    \caption{How the three cascade stages resolved the 80 test transactions. Most cases never reach the language model.}
    \label{fig:cascade-usage}
\end{figure}

\section{Discussion}
\label{sec:result-discussion}

Taken together, the results support the project's central claim: pairing a
deterministic tax engine with a multi-agent critic loop and AIS reconciliation
beats a monolithic language model on Indian tax scenarios by a wide and, more
importantly, an explainable margin --- about twenty points of tax accuracy and a
roughly three-and-a-half-fold reduction in hallucination. The explainability is
the part worth emphasising. The advantage does not come from a larger model or a
cleverer prompt; it comes from a structural decision about which sub-problems to
hand to deterministic code and which to leave to the model, and the ablation
makes that attribution precise rather than hand-waved.

It is equally important to be candid about where the system falls short, because
the gaps are real and instructive. The clearest individual failure is a case in
which a return with undeclared crypto was scored HIGH rather than the expected
CRITICAL. The cause is structural in the additive risk model: undeclared crypto
alone scores 60, which lands in the HIGH band, and escalation to CRITICAL
requires a second corroborating signal such as an explicit 194S TDS entry or a
broker CSV. When the only evidence is a bank deposit into a crypto exchange,
that second signal is absent and the score stays one band too low. The Auditor
correctly identified the crypto; the scale-up rule simply did not fire. A second
known gap is in Schedule OS recall, which sits at 87\% precisely because the
system refuses to assert the schedule without a document signal --- a deliberate
precision-for-recall trade that could be partly reversed with a low-threshold
income heuristic at a few points' cost in precision. A third is the held-out
degradation under injected noise, which is the deterministic engine faithfully
computing on slightly-wrong OCR inputs, and which points to input pre-processing
as the natural fix. The latency, finally, is the honest cost of running a 7B
model locally; it is slower than a cloud API and that is the trade the privacy
requirement demands. None of these undermine the core finding, and each has a
concrete mitigation path, which the final chapter takes up.
